\section{Competition and information}\label{sec:payoffs}

\subsection{Acquisition profits and target value}\label{sec:profits}

The auction determines both the return to preparation and the claim that informed investors trade. Consider a
prepared challenger with value \(\theta\ge p\). If it wins, it pays \(\max\{p,R\}\); if it loses, the
incumbent pays \(\theta\). The target therefore receives \(\max\{p,\min(R,\theta)\}\), and the challenger
keeps \((\theta-\max\{p,R\})_+\). A prepared challenger with \(\theta<p\) cannot bid, and the target then
receives the reserve if the incumbent's value meets it, exactly as without entry. For every value at or above
the reserve, the two expected amounts add up to the value:
\[
t(\theta)+g(\theta)=\theta,\qquad \theta\ge p.
\]

The second-price rule passes the whole value of an admissible challenger to the target and the challenger
together. An admissible value also raises proceeds to at least the reserve, so \(t_j\ge t_0\) for both
classes.

In the point-mass benchmark with the uniform incumbent, integration over \(R\) gives
\begin{equation}\label{eq:payoffs}
\begin{aligned}
t_0&=p\Bigl(1-\frac pr\Bigr),&
t_H&=\frac r2+\frac{p^2}{2r},\\
t_L&=\ell-\frac{\ell^2-p^2}{2r},&
g_H&=h-\frac r2-\frac{p^2}{2r},\\
g_L&=\frac{\ell^2-p^2}{2r},&
\DT(r)&=\frac{(r-\ell)^2}{2r}.
\end{aligned}
\end{equation}
The relevant derivatives are
\begin{equation}\label{eq:derivatives}
\DT'(r)=\frac12-\frac{\ell^2}{2r^2}>0,\qquad
g_H'(r)=-\frac12+\frac{p^2}{2r^2}<0,\qquad
g_L'(r)=-\frac{\ell^2-p^2}{2r^2}<0.
\end{equation}

Thus a stronger incumbent lowers acquisition profit at every fixed belief while it makes target proceeds more
sensitive to challenger quality. Proposition~\ref{prop:returns} extends this opposition to every continuous
incumbent distribution and every pair of classes that satisfies the class ranking~\eqref{eq:classes}.

\begin{proposition}[Competition and the two returns to information; analytical]\label{prop:returns}
Let the incumbent's value have a continuous distribution \(F\) on \([0,\bar r]\), extended by unity above
\(\bar r\). For any classes \(G_H,G_L\) with finite means,
\begin{equation}\label{eq:spread}
\begin{aligned}
\DT(F)&=pF(p)D(p^-)+\int_p^\infty[1-F(y)]\,D(y)\,dy,\\
g_j(F)&=\int_p^\infty[1-G_j(y)]\,F(y)\,dy,\qquad j\in\{H,L\}.
\end{aligned}
\end{equation}
Let \(\tilde F\le F\) be a continuous first-order stochastic strengthening of \(F\).
\begin{enumerate}[label=(\alph*),leftmargin=2em]
\item For every pair of classes, the strengthening weakly decreases \(g_H\) and \(g_L\), and hence gross
acquisition profit at every fixed belief.
\item Under the class ranking~\eqref{eq:classes}, the strengthening weakly increases the target-payoff spread.
\end{enumerate}
The decrease in \(g_j\) is strict when \(\int_p^\infty[1-G_j(y)][F(y)-\tilde F(y)]\,dy>0\). The increase in
(b) is strict when \(\int_p^\infty[F(y)-\tilde F(y)]\,D(y)\,dy>0\).
\end{proposition}

The proof works one incumbent value at a time. At a value \(R\) that meets the reserve, the target receives
\(\min\{R,\max\{\theta,p\}\}\), and the expected class difference in proceeds is then
\(\int_p^RD(y)\,dy\). Read the integrand as slices \(D(y)\,dy\) of the class value gap at each level
\(y>p\): the slice at \(y\) reaches the target when the incumbent outbids the challenger at that level and
pays its bid, and it stays with the challenger otherwise.

A stronger incumbent moves slices to the target, and the spread rises while acquisition profit falls. The
reserve adds one term: below the reserve, only an admissible challenger buys, at \(p\), and the class
difference is \(pD(p^-)\). A stronger incumbent fails the reserve less often, so that term shrinks, and it
can outweigh the moved slices when the low class is more often inadmissible. Under the class
ranking~\eqref{eq:classes}, every value meets the reserve, and the term is zero.
Appendix~\ref{app:payoffs} gives the proof.

The sum of the two lines of \eqref{eq:spread} under the class ranking~\eqref{eq:classes} is
\[
\DT(F)+g_H(F)-g_L(F)=\E_H[\max\{\theta,p\}]-\E_L[\max\{\theta,p\}].
\]
The identity holds for every incumbent distribution, and in the benchmark the right side is \(h-\ell\) at
every \(r\). A stronger incumbent leaves the value gap between the classes unchanged; it moves the gap from
the challenger's profit to target proceeds.

The spread approaches zero as \(r\) approaches \(\ell\), since the incumbent then almost never outbids a
low-value challenger. As \(r\) increases, the spread rises while gross acquisition profit falls at each
belief. Whether entry rises depends on how this change in payoffs affects the information in the price.

\subsection{What the challenger learns from the price}\label{sec:inference}

The challenger sees one price. A liquidity order in the same direction is always possible, and this bounds
what any order can reveal. Each class reaches a given nonzero order through a liquidity trader with
probability \(b\), and the informed investor adds at most \(\lambda\), so the likelihood ratio of the order
lies between \(1/(1+u)\) and \(1+u\). By Bayes' rule, the posterior after every order lies in the interval
\([m(\pi),M(\pi)]\) of \eqref{eq:bounds}, under every investor strategy, mixed or pure; when the high class
buys and the low class sells, the posteriors after a sale, no order, and a purchase are \(m(\pi)\), \(\pi\),
and \(M(\pi)\). The challenger's belief at a price averages the posteriors of the orders with that price, so
it obeys the same bounds. A preparation decision that needs a belief above the top posterior therefore
follows no price: when \(g_H\ge g_L\), gross profit does not fall with the belief, so once gross profit at
the top posterior falls below \(c_H\), no order recruits a high-cost challenger, in or out of equilibrium.
This is the ceiling of the model.

When \(g_H>g_L\), gross profit covers the high cost exactly when the belief reaches the preparation
threshold \(\tau=(c_H-g_L)/(g_H-g_L)\): \(B(\mu)\ge c_H\) if and only if \(\mu\ge\tau\). For the uniform
incumbent I write \(\tau(r)\). A zero-cost challenger prepares after every price, so entry after every order
is at least \(\rho\), and competitive pricing gives \eqref{eq:price}. When the classes differ in target
proceeds, the price rises strictly with the posterior. Suppose two orders have one price. Entry after them
is the same, because entry depends on the order only through the price, and at least \(\rho\); the price
formula then gives the two orders the same posterior. The challenger therefore reads the market makers'
posterior from the price, and its gross profit at the price \(P(x)\) is \(B(\mu(x))\), so high-cost
preparation follows exactly the orders whose posterior reaches \(\tau\).
Appendix~\ref{app:beliefs} gives the formal statements and their proofs.

The zero-cost challenger gives each step its floor: entry stays at least \(\rho\), so the price reveals the
posterior, and the investor's return after every order stays at least a fixed share of the spread, the
trading bound of Section~\ref{sec:informed}. Its preparation does not respond to the price, so information
can only add entry.

\subsection{Informed trading}\label{sec:informed}

Rational pricing removes the anticipated increase in proceeds from the investor's advantage. What remains is
the gross return of a purchase in class \(H\) and of a sale in class \(L\),
\[
K_H=e(1)\,\DT\,[1-\mu(1)],\qquad K_L=e(-1)\,\DT\,\mu(-1),\qquad \rho\,\kappa(\pi)\,\DT\le K_H,\,K_L\le\DT,
\]
with \(\kappa(\pi)=\min\{m(\pi),1-M(\pi)\}\). A wrong-signed order earns at most \(-k\), so it never pays.
Class \(H\) buys when \(K_H\) exceeds \(k\) and abstains when it falls short, and class \(L\) compares
\(K_L\) with \(k\) in the same way.

Market makers price the entry response, but a liquidity order in the same direction is always possible, so
some uncertainty about the class survives every order. The investor's gross return is the entry probability
times the spread times that residual uncertainty. The lower bound uses the entry floor \(\rho\) and the
posterior bounds, and the upper bound is the spread. Both bounds hold under every candidate profile, mixed or
pure, and their comparison with \(k\) decides when information survives in equilibrium.

Since \(1-M(\pi)=m(1-\pi)\), the bound takes the form \(\kappa(\pi)=m(\min\{\pi,1-\pi\})\). It is largest
at the prior \(1/2\), where it equals \(m(1/2)=1/(2+u)\), and smaller at every other prior. The class that
the prior favors binds, because the price already leans toward it.

The no-entry proceeds \(t_0\) cancel from each return, because they move the price and the terminal payoff
together; the terms that matter are the entry gains \(t_H-t_0\) and \(t_L-t_0\), which are nonnegative.

Each return depends only on the posterior and entry after its own order, which splits the investor's problem:
entry after the zero order enters no payoff, so each class compares one gross return with the trading cost,
class \(H\) the return of a purchase and class \(L\) the return of a sale.
Appendix~\ref{app:beliefs} gives the formal statement and proof of the trading bound, and
Appendix~\ref{app:coexistence} uses this split to list every equilibrium along the benchmark path.
